% !TEX root = ../notes.tex

\documentclass[../notes.tex]{subfiles}

\begin{document}

\section{April 13}
Welcome back. We continue discussing automorphic forms for function fields.

\subsection{The Central Character}
% Let's begin by discussing how to extend our definition of automorphic form for a reductive group.
% \begin{definition}
% 	Fix a reductive group $G$ over a function field $F$ with center $Z$. Given a continuous character $\omega\colon Z(F)\backslash Z(\AA_F)\to\CC^\times$, we define $\mc A(\omega)$ to be the subspace of $Z(\AA_F)$-eigenvectors of $C^\infty(G(F)\backslash G(\AA_F))$ with eigenvalue $\omega$.
% \end{definition}
We did not bother showing that the space of cusp forms was finite-dimensional over number fields. Today, we will content ourselves with sketching the result over function fields.
\begin{theorem} \label{thm:cusp-compact-support}
	Fix a semisimple group $G$ over a function field $F$, and let $K\subseteq G(\AA_F)$ be a compact open subgroup.
	\begin{listalph}
		\item The space $\mc A_{\mathrm{cusp}}(G,K)$ is finite-dimensional.
		\item There is a finite subset $U_K\subseteq G(F)\backslash G(\AA_F)/K$ such that all cups forms in $\mc A_{\mathrm{cusp}}(K)$ vanish outside $U_K$.
	\end{listalph}
\end{theorem}
\begin{remark}
	It follows that any cusp form for $G$ has support with compact image in $G(F)\backslash G(\AA_F)$.
\end{remark}
\begin{remark}
	Over number fields, the analog of this result is that cusp forms have rapid decay at the cusps; this strengthens the ``moderate growth'' at the cusps condition required for general automorphic forms. For example, for $\op{GL}_{2,\QQ}$, this implies that the cuspidal holomorphic modular forms $f(x+iy)$ are $o\left(y^{-N}\right)$ for all $N$. We can see this directly: using the $q$-expansion, one actually sees that $f(x+iy)$ is $O(\exp(-y))$.
\end{remark}
We will sketch this result only for $G=\op{GL}_n$ and $K=K_0$. For this, we should figure out what the correct statement is when $G$ is not semisimple. The right thing to do is to fix the central character, which roughly speaking kills the torus.
\begin{definition}[central character]
	Fix a reductive group $G$ over a function field $F$ with center $Z$. Given a continuous character $\omega\colon Z(F)\backslash Z(\AA_F)\to\CC^\times$, we see that $f\in\mc A(G)$ has \textit{central character} $\omega$ if and only if $f$ is a $Z(\AA_F)$-eigenvector with eigenvalue $\omega$. We let $\mc A(\omega)$ denote the subspace of automorphic forms with central character $\omega$.
\end{definition}
\begin{remark}
	Morally, fixing the central character passes the automorphic form from $G(F)\backslash G(\AA_F)$ to the quotient $G(F)\backslash G(\AA_F)/Z(\AA_F)$. The former has no chance of being finite volume when $Z$ has positive dimension because the norm map surjects onto $q^\ZZ$, but it turns out that the latter has finite volume.
\end{remark}
\begin{remark}
	Because $Z$ is central, we see that
	\[\mc A(G)=\bigsqcup_\omega\mc A(\omega).\]
\end{remark}
\begin{example}
	Take $G=\op{GL}_n(F)$ with level $K_0=\prod_{x\in\left|X\right|}\op{GL}_n(\OO_x)$. Using \Cref{exe:double-quotient-as-groupoid}, we see that an automorphic form $f$ with central character $\omega$ is required to have
	\[f(\mc E\otimes\mc L)=\omega(\mc L)f(\mc E).\]
	Thus, by letting $\mc L$ vary over various degrees, we see that $f$ is determined by its values on vector bundles of degrees in $\{0,1,\ldots,n-1\}$.
\end{example}
The analog of \Cref{thm:cusp-compact-support} which we will show is the following.
\begin{notation}
	Fix a curve $X$. Then we let $\op{Vect}_n^d(X)$ be the moduli stack of vector bundles of rank $n$ and degree $d$.
\end{notation}
\begin{remark}
	There is an equivalence
	\[\op{Vect}_n(X)=\bigsqcup_{d\in\ZZ}\op{Vect}_n(X).\]
\end{remark}
\begin{proposition} \label{prop:function-field-cuspidal-fd}
	Fix a function field $F=\FF_q(X)$, and set $K_0\coloneqq\prod_{x\in\left|X\right|}\op{GL}_n(\OO_x)$. Then there is a finite subset
	\[U\subseteq\bigsqcup_{i=0}^{n-1}\op{Vect}_n^i(X)\]
	such that $f$ vanishes outside $\op{Pic}(X)\cdot U\subseteq\op{Vect}_n(X)$.
\end{proposition}

\subsection{Cusp Forms for Function Fields}
Let's start by explaining what it means for an automorphic form to be cuspidal.
\begin{lemma}
	Fix a function field $F=\FF_q(X)$, and set $K_0\coloneqq\prod_{x\in\left|X\right|}\op{GL}_n(\OO_x)$. Then an automorphic form $f$ is cuspidal if and only if
	\[\sum_{\mc E\in\op{Ext}^1(\mc E'',\mc E')}f(\mc E)=0\]
	for all vector bundles $\mc E'$ and $\mc E''$ of ranks $m$ and $n-m$, respectively.
\end{lemma}
\begin{proof}
	It is enough to check cuspidality on maximal parabolic subgroups $P\subseteq\op{GL}_n(F)$. Once $P$ contains the standard Borel of upper-triangular matrices, $P$ looks like block-diagonal matrices $\begin{bsmallmatrix}
		A & B \\ 0 & D
	\end{bsmallmatrix}$, where $A$ is $m\times m$ and $D$ is $(n-m)\times(n-m)$. Now, let $P=MN$ be the Levi decomposition, and we are interested in checking when
	\[\op{CT}_P(g)\coloneqq\int_{N(F)\backslash N(\AA_F)}f(ng)\,dn\]
	vanishes for all $g\in G(\AA_F)$. Well, the Iwasawa decomposition at each place turns out to induce an Iwasawa decomposition
	\[G(\AA_F)=N(\AA_F)M(\AA_F)K_0,\]
	so the function $\op{CT}_P$ descends to the quotient $M(F)\backslash M(\AA_F)/K_{M,0}$, where $K_{M,0}\coloneqq K_0\cap M(\AA_F)$. (Here, $M(F)$ comes out on the left because $M$ normalizes $N$.) To compute $\op{CT}_P$, we note that $M=\mathrm{GL}_m\times\mathrm{GL}_{n-m}$, so the source $M(F)\backslash M(\AA_F)/K_{M,0}$ simply parameterizes pairs $(\mc E',\mc E'')$ of vector bundles of ranks $m$ and $n-m$, respectively. Then we claim that
	\[\op{CT}_P(\mc E',\mc E'')=\sum_{\mc E\in\op{Ext}^1(\mc E'',\mc E')}f(\mc E)\cdot\frac1{\#\op{Aut}(\mc E)},\]
	where $f(\mc E)$ refers to the value of $f$ on the underlying vector bundle of $\mc E$. Indeed, for $g\in M(\AA_F)$, we see that the choice of $n\in N(\AA_F)$ grants us some filling in for a short exact sequence
	\[0\to\mc E'\to\mc E\to\mc E''\to0,\]
	and then the integral simply sums the value of $f(\mc E)$ as the short exact sequence varies.
	
	Lastly, we have to remove the automorphism group from our sum. This follows because $\op{Aut}(\mc E)$ actually has constant size: an automorphism of the short exact sequence is an automorphism $\alpha$ of the underlying vector bundle $\mc E$ which restricts to the identity on the sub-bundle $\mc E'$ and the quotient $\mc E''$. But then $\alpha-\id_{\mc E}$ gives rise to a projection $\mc E''\to\mc E'$, which splits and thus trivializes the sequence, so $\alpha-\id_{\mc E}$ is then parameterized by $\op{Hom}(\mc E'',\mc E')$.
\end{proof}

\subsection{Reduction Theory for Function Fields}
Next, we need some reduction theory. Over number fields, reduction theory for $\op{GL}_n$ arises from choosing some basis of shortest vectors of the associated metrized lattices. The analog of a vector in a vector bundle $\mc E$ is a line sub-bundle $\mc L\subseteq\mc E$, and the analog of being shortest is having large degree.\footnote{For example, one can expect to find line bundles of smaller degree by appropriately twisting, so it is only ever surprising to have sub-bundles of large degree.} This motivates the following discussion.
\begin{definition}[slope]
	Fix a vector bundle $\mc E$ on a curve $X$. Then the \textit{slope} is
	\[\mu(\mc E)\coloneqq\frac{\deg\mc E}{\op{rank}\mc E}.\]
\end{definition}
\begin{definition}[semistable]
	Fix a vector bundle $\mc E$ on a curve $X$. Then $\mc E$ is \textit{semistable} if and only if all nonzero sub-bundles $\mc F\subseteq\mc E$ have $\mu(F)\le\mu(E)$.
\end{definition}
\begin{example}
	Suppose that $\mc E$ is a direct sum of two line bundles $\mc L_1$ and $\mc L_2$. If $\deg\mc L_1>\deg\mc L_2$, then $\mc E$ is not semistable (which we can see by inspection). However, if the degrees are equal, then $\mc E$ is semistable. Indeed, for any $\mc L\subseteq\mc L_1\oplus\mc L_2$, one can show that
	\[\deg\mc L\le\max\{\deg\mc L_1,\deg\mc L_2\}.\]
\end{example}
\begin{theorem}[Harder--Narasimhan]
	Fix a vector bundle $\mc E$ on a curve $X$. Then $\mc E$ admits a unique filtration
	\[0=\mc F_0\subsetneq\mc F_1\subsetneq\mc F_2\subsetneq\cdots\subsetneq \mc F_r=\mc E\]
	for which $\mc F_i/\mc F_{i-1}$ is semistable for each $i$, and $\{\mu(\mc F_i/\mc F_{i-1})\}_i$ is decreasing.
\end{theorem}
\begin{proof}[Sketch]
	The main point is to choose $\mc F_1$, from which the rest of the filtration can be found inductively by passing to the quotient. To this end, we simply choose $\mc F_1$ to have maximal slope, and among those sub-bundles with maximal slope, we choose $\mc F_1$ to have maximal rank. It is nontrivial to verify uniqueness.
\end{proof}
This filtration allows us to define an interesting invariant.
\begin{definition}[polygon]
	Fix a vector bundle $\mc E$ on a curve $X$ with filtration. Then its \textit{polygon} $\nu$ is the line connecting $(0,0)$ to the points $(\op{rank}\mc F_i,\deg\mc F_i)$ in order. We let $\op{Vect}_n^\nu(X)$ denote the collection of vector bundles of rank $n$ and polygon $\nu$.
\end{definition}
\begin{remark}
	The construction of the filtration implies that the polygon is concave down.
\end{remark}
\begin{remark}
	It turns out that each $\op{Vect}^\nu_n(X)$ is finite. Roughly speaking, this follows because there are only finitely many semistable bundles of given rank.
\end{remark}
Here is our argument.
\begin{proof}[Proof of \Cref{prop:function-field-cuspidal-fd}]
	It suffices to show that there is a polygon $\nu_0$ with endpoints $(n,d)$ such that all cusp forms on $\op{Vect}_n^d(X)$ must vanish on $\op{Vect}^\nu_n(X)$ when $\nu\not\le\nu_0$. Namely, we recover the result by letting $d$ vary over our finitely many degrees, and the finite set is simply given by those vector bundles with polygon bounded above by $\nu_0$.

	For simplicity, we now take $G=\op{GL}_2$ and $d=0$ so that the choice of polygon $\nu_0$ amounts to a single choice of point $(1,e_0)$; the rest of the polygon comes from the points $(0,0)$ and $(2,0)$. Now, via the filtration, we see that $\op{Vect}_2^{(1,e)}(X)$ is parameterized by short exact sequences
	\[0\to\mc E'\to\mc E\to\mc E''\to0\]
	where $\mc E'$ has degree $e$, $\mc E$ has degree zero, and $\mc E''$ has degree $-e$. Now, for our cuspidality check, we should fix line bundles $\mc E'$ and $\mc E''$ of degrees $e$ and $-e$, and we are given that
	\[\sum_{\mc E\in\op{Ext}^1(\mc E'',\mc E')}f(\mc E)=0.\]
	Now, for $e$ sufficiently large, we note that the short exact sequence splits: indeed, $\op{Ext}^1(\mc E'',\mc E')$ is some cohomology group in $\mc E''\otimes(\mc E')^\lor$, so high degrees induces vanishing by Serre vanishing. Thus, this cuspidality condition directly implies that $f(\mc E)=0$ for $e$ large enough. This completes the argument in this case.
\end{proof}
\begin{remark}
	For $\op{GL}_n$ for higher $n$, one needs a more complicated induction.
\end{remark}

\subsection{The Cuspidal Subspace}
Over number fields, the cuspidal subspace decomposes as follows.
\begin{proposition}
	Fix a semisimple group $G$ over a number field $F$. Let $L^2_{\mathrm{cusp}}\subseteq L^2(G(F)\backslash G(\AA_F))$ be the closure of the subspace of cusp forms. Then $L^2_{\mathrm{cusp}}$ is a Hilbert space direct sum of irreducible representations of $G(\AA_F)$, and they occur with finite multiplicity.
\end{proposition}
\begin{proof}
	We will show this in the special case where $G(F)\backslash G(\AA_F)$ is compact. In this case, it turns out that all automorphic forms are automatically cusp forms; roughly speaking, this holds because there are no proper parabolic subgroups because $G$ is anisotropic. (Parabolic subgroups can be defined as ``attracting subsets'' of one-parameter subgroups $\mathbb G_m\to G$.)

	We start by choosing an approximation to identity $\{f_n\}$ in $C_c(G(\AA_F),\RR)$. By replacing $f_n$ with $g\mapsto\frac12\left(f_n(g)+f_n(g^{-1})\right)$, we may assume that $f_n(g)=f_n\left(g^{-1}\right)$ for all $g$. Now, note that the induced action $R(f_n)$ on $C_c^\infty(G(\AA_F),\RR)$ is (by definition) given by a kernel. Thus, they are Hilbert--Schmidt operators and thus compact, and they are self-adjoint by the symmetry property, so we use the spectral theorem of compact operators. Namely, each $R(f_n)$ diagonalizes into a spectrum $\{\lambda_n\}_n$ as
	\[L^2(G(F)\backslash G(\AA_F))[0]\oplus\bigoplus_nL^2(G(F)\backslash G(\AA_F))[\lambda_n],\]
	and the sequence $\{\lambda_n\}_n$ (with multiplicities) is discrete, except possibly at zero.
	
	Now, every vector in $L^2(G(F)\backslash G(\AA_F))$ has some nontrivial action by some $R(f_n)$ because $\delta_1$ acts nontrivially (indeed, it fixes every vector). Thus, we decompose into nonzero eigenspaces, and it only remains to show that our spaces occur with finite multiplicity. Well, looping over all $f_n$s, one can show a given invariant subspace occurs with finite multiplicity, and the decomposition follows because the ma
\end{proof}
\begin{remark}
	This is the same argument as the proof of the Peter--Weyl theorem.
\end{remark}

\end{document}